\documentclass[11pt,a4paper]{article}
\usepackage[utf8]{inputenc}
\usepackage[T1]{fontenc}
\usepackage{amsmath,amssymb,amsthm}
\usepackage{booktabs}
\usepackage{geometry}
\usepackage{xcolor}
\usepackage{hyperref}
\usepackage{tcolorbox}
\usepackage{enumitem}

\geometry{margin=1in}
\definecolor{bonboblue}{RGB}{0,90,160}

\tcbuselibrary{breakable}
\newtcolorbox{formulabox}[1][]{
  colback=blue!3!white,
  colframe=bonboblue!80!black,
  fonttitle=\bfseries,
  breakable,
  #1
}

\title{\textbf{Slippage Minimization: Mathematical Formulas} \\[0.3em]
\large BonBoExtend Execution Engine — Formal Specification}
\author{BonBo Execution Engine v1.0}
\date{\today}

\begin{document}
\maketitle
\tableofcontents
\newpage

%════════════════════════════════════════════════════════════════
\section{Square-Root Market Impact Model}
%════════════════════════════════════════════════════════════════

The core impact model follows the empirically-validated \textbf{square-root law}:

\begin{formulabox}[title={\S1.1 — Temporary Impact (Square-Root Law)}]
Let $Q$ be the order notional (USD), $V$ the average daily volume, $\sigma$ the daily volatility, and $\eta$ the temporary impact coefficient. Then:
\begin{equation}
  \mathcal{I}_{\text{temp}} = \eta \cdot \sigma \cdot \sqrt{\frac{Q}{V}} \times 10{,}000 \quad \text{(bps)}
\end{equation}
where $\sqrt{Q/V}$ is the \textbf{participation rate} square root. This is the dominant component of market impact for large orders.
\end{formulabox}

\begin{formulabox}[title={\S1.2 — Permanent Impact (Linear Model)}]
The non-revertible permanent impact is linear in participation rate, with coefficient $\gamma$:
\begin{equation}
  \mathcal{I}_{\text{perm}} = \gamma \cdot \frac{Q}{V} \times 10{,}000 \quad \text{(bps)}
\end{equation}
Typically $\gamma \ll \eta$, so permanent impact is much smaller than temporary.
\end{formulabox}

\begin{formulabox}[title={\S1.3 — Total Market Impact}]
\begin{equation}
  \mathcal{I}_{\text{total}} = \underbrace{\eta \cdot \sigma \cdot \sqrt{\frac{Q}{V}}}_{\text{temporary}} + \underbrace{\gamma \cdot \frac{Q}{V}}_{\text{permanent}} \quad \text{(bps)}
\end{equation}
\end{formulabox}

\begin{formulabox}[title={\S1.4 — Risk-Adjusted Cost (CVaR Penalty)}]
With risk aversion parameter $\lambda$:
\begin{equation}
  \mathcal{C}_{\text{risk}} = \mathcal{I}_{\text{total}} + f_{\text{fee}} + \underbrace{\lambda \cdot \sigma \cdot \sqrt{\frac{Q}{V}} \times 5{,}000}_{\text{CVaR penalty}} \quad \text{(bps)}
\end{equation}
where $f_{\text{fee}}$ is the taker fee rate in bps.
\end{formulabox}

\begin{formulabox}[title={\S1.5 — Calibrated Parameters}]
\begin{center}
\begin{tabular}{lcccc}
\toprule
\textbf{Symbol} & $\sigma$ & $\eta$ & $\gamma$ & $V_{\text{daily}}$ \\
\midrule
BTCUSDT  & 0.025 & 0.15 & 0.01 & \$5B  \\
ETHUSDT  & 0.030 & 0.20 & 0.015 & \$2B \\
SOLUSDT  & 0.050 & 0.30 & 0.025 & \$500M \\
SEIUSDT  & 0.060 & 0.40 & 0.035 & \$100M \\
\bottomrule
\end{tabular}
\end{center}
\end{formulabox}


%════════════════════════════════════════════════════════════════
\section{Almgren-Chriss Optimal Execution}
%════════════════════════════════════════════════════════════════

\begin{formulabox}[title={\S2.1 — Urgency Parameter}]
The urgency (or ``trading rate'') parameter $\kappa$ balances trading cost against timing risk:
\begin{equation}
  \kappa = \sqrt{\frac{\lambda \cdot \sigma^2}{\eta}}
\end{equation}
where $\lambda$ is risk aversion, $\sigma$ is volatility, and $\eta$ is the temporary impact coefficient.
\end{formulabox}

\begin{formulabox}[title={\S2.2 — Optimal Execution Time}]
\begin{equation}
  T^{*} = \sqrt{\frac{3\eta}{\lambda \cdot \sigma^2}} \cdot \left(1 + 100 \cdot \frac{Q}{V}\right) \quad \text{(hours)}
\end{equation}
This is the classical Almgren-Chiss result scaled by participation rate for crypto markets. Clamped to $[1\text{min},\; 8\text{h}]$.
\end{formulabox}

\begin{formulabox}[title={\S2.3 — Optimal Trading Trajectory}]
The fraction of total quantity remaining at time $t$ in the execution window $[0, T]$:
\begin{equation}
  x(t) = \frac{\sinh\!\bigl(\kappa(T - t)\bigr)}{\sinh(\kappa T)}
\end{equation}
This gives a \textbf{front-loaded} trajectory: faster execution early, slower later.
\end{formulabox}

\begin{formulabox}[title={\S2.4 — Slice Fractions}]
For $N$ slices over execution time $T$, the $i$-th slice fraction is:
\begin{equation}
  f_i = x(t_i) - x(t_{i+1}) = \frac{\sinh\!\bigl(\kappa(T - t_i)\bigr) - \sinh\!\bigl(\kappa(T - t_{i+1})\bigr)}{\sinh(\kappa T)}
\end{equation}
where $t_i = \frac{i}{N} \cdot T$. Each slice quantity is $q_i = f_i \cdot Q_{\text{total}}$.

\vspace{0.5em}
\textbf{Properties:}
\begin{itemize}[nosep]
  \item $f_0 > f_1 > \cdots > f_{N-1}$: front-loaded (largest slices first)
  \item $\sum_{i=0}^{N-1} f_i = 1$: complete execution
  \item As $\kappa \to 0$: $f_i \to 1/N$ (uniform TWAP-like schedule)
  \item As $\kappa \to \infty$: $f_0 \to 1$ (immediate execution)
\end{itemize}
\end{formulabox}


%════════════════════════════════════════════════════════════════
\section{Implementation Shortfall Decomposition}
%════════════════════════════════════════════════════════════════

\begin{formulabox}[title={\S3.1 — IS Decomposition}]
\begin{equation}
  \text{IS} = \underbrace{\eta \sigma \sqrt{\frac{Q}{V}}}_{\text{temporary impact}} + \underbrace{\gamma \frac{Q}{V}}_{\text{permanent impact}} + \underbrace{\sigma \sqrt{\frac{T}{24}}}_{\text{timing risk}} + \underbrace{f_{\text{fee}}}_{\text{fees}}
\end{equation}
All terms in bps.
\end{formulabox}

\begin{formulabox}[title={\S3.2 — IS at Confidence Intervals}]
\begin{align}
  \text{IS}_{95\%} &= \mathcal{I}_{\text{expected}} + 1.645 \cdot \sigma\sqrt{T/24} \times 10{,}000 \\
  \text{IS}_{99\%} &= \mathcal{I}_{\text{expected}} + 2.326 \cdot \sigma\sqrt{T/24} \times 10{,}000
\end{align}
where $\sigma\sqrt{T/24}$ is the timing risk (volatility exposure during execution).
\end{formulabox}


%════════════════════════════════════════════════════════════════
\section{Transient Impact Model}
%════════════════════════════════════════════════════════════════

\begin{formulabox}[title={\S4.1 — Exponential Decay Kernel}]
Each prior trade at time $\tau_j$ with execution rate $r_j$ (USD/sec) contributes to current impact via an exponential decay kernel:
\begin{equation}
  \mathcal{I}_{\text{transient}}(t) = \eta \times 10^{-4} \sum_{j=1}^{M} r_j \cdot e^{-(t - \tau_j)/\tau_d}
\end{equation}
where $\tau_d$ is the decay timescale (in seconds), calibrated per symbol.
\end{formulabox}

\begin{formulabox}[title={\S4.2 — Forward Impact Projection}]
Estimated remaining impact at future time $t + \Delta t$:
\begin{equation}
  \hat{\mathcal{I}}(t + \Delta t) = \eta \times 10^{-4} \sum_{j=1}^{M} r_j \cdot e^{-(t + \Delta t - \tau_j)/\tau_d}
\end{equation}
Used to predict how much impact ``decays away'' before the next slice.
\end{formulabox}

\begin{formulabox}[title={\S4.3 — Pruning Rule}]
Trades older than $5\tau_d$ are pruned from history:
\begin{equation}
  \text{cutoff} = t - 5\tau_d
\end{equation}
Since $e^{-5} \approx 0.007$, these contribute $<0.7\%$ of total transient impact.
\end{formulabox}


%════════════════════════════════════════════════════════════════
\section{Optimal Slice Sizing (Depth-Based)}
%════════════════════════════════════════════════════════════════

\begin{formulabox}[title={\S5.1 — Depth Walk (Cumulative VWAP)}]
Walking $L$ levels of the order book from the touch:
\begin{equation}
  \text{cum\_qty}_i = \sum_{j=0}^{i} q_j, \qquad
  \text{cum\_cost}_i = \sum_{j=0}^{i} q_j \cdot p_j
\end{equation}
\begin{equation}
  \text{VWAP}_i = \frac{\text{cum\_cost}_i}{\text{cum\_qty}_i}, \qquad
  \text{impact}_i = \frac{|\text{VWAP}_i - P_{\text{mid}}|}{P_{\text{mid}}} \times 10{,}000
\end{equation}
\end{formulabox}

\begin{formulabox}[title={\S5.2 — Slippage Budget Constraint}]
The optimal slice quantity $Q^{*}$ is the maximum $\text{cum\_qty}$ where impact $\leq \beta$ (max bps budget):
\begin{equation}
  Q^{*} = \max\left\{\text{cum\_qty}_i \;\Bigg|\; \text{impact}_i \leq \beta\right\}
\end{equation}
\end{formulabox}

\begin{formulabox}[title={\S5.3 — Level Interpolation}]
When the budget $\beta$ is exceeded mid-level, linear interpolation finds the exact quantity:
\begin{equation}
  Q^{*} = \text{cum\_qty}_{i-1} + q_i \cdot \frac{\beta - \text{impact}_{i-1}}{\text{impact}_i - \text{impact}_{i-1}}
\end{equation}
This avoids over-conservative rounding down to the previous level.
\end{formulabox}

\begin{formulabox}[title={\S5.4 — Participation Rate Cap}]
\begin{equation}
  Q_{\text{cap}} = \rho_{\max} \cdot \sum_{j=0}^{L} q_j
\end{equation}
where $\rho_{\max}$ is the maximum participation rate (e.g.\ 10\%). The final slice respects:
$Q_{\text{slice}} = \min(Q^{*},\; Q_{\text{cap}},\; Q_{\text{remaining}})$.
\end{formulabox}

\begin{formulabox}[title={\S5.5 — Transient Consumption Adjustment}]
Prior slices leave a decaying ``footprint'' on the book. The effective consumption:
\begin{equation}
  C_{\text{transient}} = \sum_{k=0}^{K-1} q_{k}^{\text{slice}} \cdot \delta^k
\end{equation}
where $\delta$ is the decay factor (e.g.\ 0.7) and $k=0$ is the most recent slice.
The available depth is reduced: $Q_{\text{avail}} = Q^{*} - C_{\text{transient}}$.
\end{formulabox}

\begin{formulabox}[title={\S5.6 — OFI Size Multiplier}]
The Order Flow Imbalance signal adjusts the slice size:
\begin{equation}
  Q_{\text{adjusted}} = Q_{\text{avail}} \times \mu_{\text{OFI}}
\end{equation}
where the multiplier $\mu_{\text{OFI}}$ depends on signal strength:
\begin{center}
\begin{tabular}{lcc}
\toprule
\textbf{Signal} & \textbf{Favorable} & $\mu_{\text{OFI}}$ \\
\midrule
StrongBuy / StrongSell & Yes & 1.3 (boost) \\
Buy / Sell             & Yes & 1.15 (mild boost) \\
Neutral                & —   & 1.0 (unchanged) \\
Buy / Sell             & No  & 0.85 (mild cut) \\
StrongBuy / StrongSell & No  & 0.7 (penalty) \\
\bottomrule
\end{tabular}
\end{center}
``Favorable'' means the OFI direction aligns with our order side.
\end{formulabox}


%════════════════════════════════════════════════════════════════
\section{Order Flow Imbalance (OFI)}
%════════════════════════════════════════════════════════════════

\begin{formulabox}[title={\S6.1 — Volume Imbalance}]
\begin{equation}
  \phi = \frac{V_{\text{bid}}(L) - V_{\text{ask}}(L)}{V_{\text{bid}}(L) + V_{\text{ask}}(L)}
\end{equation}
where $V_{\text{bid}}(L)$ and $V_{\text{ask}}(L)$ are the total quantities in the top $L$ levels of each side. Range: $\phi \in [-1, +1]$.
\end{formulabox}

\begin{formulabox}[title={\S6.2 — Depth Skew}]
\begin{equation}
  \psi_{\text{depth}} = \frac{V_{\text{bid}} - V_{\text{ask}}}{V_{\text{bid}} + V_{\text{ask}}}
\end{equation}
Same formula as imbalance but computed over the full book depth. Detects deeper liquidity asymmetry.
\end{formulabox}

\begin{formulabox}[title={\S6.3 — Wall Detection}]
A wall exists at level $i$ if:
\begin{equation}
  w_i = \frac{q_i}{\displaystyle\sum_{j=0}^{L} q_j} > \theta_{\text{wall}}
\end{equation}
where $\theta_{\text{wall}}$ is the wall threshold (e.g.\ 0.35 = 35\% of total depth in a single level).
\end{formulabox}

\begin{formulabox}[title={\S6.4 — Signal Classification}]
\begin{center}
\begin{tabular}{lll}
\toprule
\textbf{Condition} & \textbf{Signal} & \textbf{Action} \\
\midrule
$\phi \geq 0.58$ & StrongBuy  & Aggressive (1.5$\times$ size) \\
$0.45 \leq \phi < 0.58$ & Buy     & Passive (1.0$\times$) \\
$\phi \approx 0.5$       & Neutral & Passive (0.7$\times$) \\
$0.42 < \phi \leq 0.45$ & Sell     & Passive (1.0$\times$) \\
$\phi < 0.42$            & StrongSell & Defensive (0.5$\times$) \\
\bottomrule
\end{tabular}
\end{center}
For SELL-side orders, the conditions are mirrored.
\end{formulabox}


%════════════════════════════════════════════════════════════════
\section{Dynamic Spread Threshold}
%════════════════════════════════════════════════════════════════

\begin{formulabox}[title={\S7.1 — Spread Measurement}]
\begin{equation}
  s = \frac{P_{\text{ask}} - P_{\text{bid}}}{\frac{1}{2}(P_{\text{ask}} + P_{\text{bid}})} \times 10{,}000 \quad \text{(bps)}
\end{equation}
\end{formulabox}

\begin{formulabox}[title={\S7.2 — Volatility-Aware Dynamic Threshold}]
Given a history $\{s_1, s_2, \ldots, s_N\}$ of recent spread readings:
\begin{align}
  \bar{s} &= \frac{1}{N}\sum_{i=1}^{N} s_i \qquad \text{(mean spread)} \\
  \sigma_s &= \sqrt{\frac{1}{N-1}\sum_{i=1}^{N}(s_i - \bar{s})^2} \qquad \text{(spread volatility)}
\end{align}
The dynamic threshold for a fixed parameter $\theta_{\text{fixed}}$:
\begin{equation}
  \theta_{\text{dyn}} = \theta_{\text{fixed}} \times \left(1 + \nu \cdot \frac{\sigma_s}{s_0}\right)
\end{equation}
where $\nu$ is the volatility multiplier and $s_0$ is the base (typical) spread.
\end{formulabox}

\begin{formulabox}[title={\S7.3 — Route Decision}]
Given current spread $s$, dynamic flash threshold $\theta_F$, market threshold $\theta_M$, and abort threshold $\theta_A$:
\begin{equation}
  \text{route}(s) = \begin{cases}
    \text{Flash Limit}  & \text{if } s \leq \theta_F \\
    \text{Adaptive Limit} & \text{if } \theta_F < s \leq \theta_M \\
    \text{Market}       & \text{if } \theta_M < s \leq \theta_A \\
    \text{Hold}         & \text{if } s > \theta_A
  \end{cases}
\end{equation}
\end{formulabox}

\begin{formulabox}[title={\S7.4 — Flash Limit Savings}]
When a flash limit fills at the touch price instead of market:
\begin{equation}
  \text{savings} = s \quad \text{(entire spread, in bps)}
\end{equation}
For BTCUSDT with typical 2bps spread and 1000 BTC order (\$100M):
\begin{equation}
  \text{savings} = 2 \text{bps} \times \$100\text{M} = \$20{,}000
\end{equation}
\end{formulabox}


%════════════════════════════════════════════════════════════════
\section{Smart-Market Book Classification}
%════════════════════════════════════════════════════════════════

\begin{formulabox}[title={\S8.1 — Depth-Weighted Mid Price}]
\begin{equation}
  P_{\text{mid}}^{*} = \frac{\displaystyle\sum_{i=0}^{L} p_i^{\text{bid}} \cdot q_i^{\text{bid}} + \sum_{i=0}^{L} p_i^{\text{ask}} \cdot q_i^{\text{ask}}}{\displaystyle\sum_{i=0}^{L} q_i^{\text{bid}} + \sum_{i=0}^{L} q_i^{\text{ask}}}
\end{equation}
Volume-weighted center of gravity of the book. More robust than simple mid when depth is asymmetric.
\end{formulabox}

\begin{formulabox}[title={\S8.2 — Effective Quantity}]
Size adjustment based on book state:
\begin{equation}
  Q_{\text{eff}} = Q_{\text{planned}} \times \alpha_{\text{state}}
\end{equation}
\begin{center}
\begin{tabular}{lc}
\toprule
\textbf{Book State} & $\alpha_{\text{state}}$ \\
\midrule
Aggressive  & $1.0$--$1.5$ \\
Passive     & $1.0$ \\
Defensive   & $0.7$ \\
WallBlocking & $0.5$ \\
WideSpread  & $0.5$ \\
\bottomrule
\end{tabular}
\end{center}
\end{formulabox}

\begin{formulabox}[title={\S8.3 — Limit Price with OFI Offset}]
\begin{equation}
  P_{\text{limit}} = P_{\text{mid}}^{*} + \text{sign}(\text{side}) \cdot \Delta_{\text{OFI}}
\end{equation}
where the offset $\Delta_{\text{OFI}}$ (in bps) is:
\begin{center}
\begin{tabular}{lc}
\toprule
\textbf{OFI State} & $\Delta_{\text{OFI}}$ (bps) \\
\midrule
Aggressive & 0 (join touch) \\
Passive    & 1 (inside spread) \\
Defensive  & 2 (deep inside) \\
\bottomrule
\end{tabular}
\end{center}
\end{formulabox}


%════════════════════════════════════════════════════════════════
\section{Order Book Slippage Estimation}
%════════════════════════════════════════════════════════════════

\begin{formulabox}[title={\S9.1 — Buy Slippage Walk}]
For a buy order of quantity $Q$, walking ask levels:
\begin{equation}
  \text{VWAP}(Q) = \frac{\displaystyle\sum_{i=0}^{n-1} p_i \cdot q_i + p_n \cdot (Q - Q_{n-1})}{Q}
\end{equation}
where $Q_{n-1} = \sum_{i=0}^{n-1} q_i < Q$ and $Q_n = Q_{n-1} + q_n \geq Q$.
\end{formulabox}

\begin{formulabox}[title={\S9.2 — Slippage Estimate}]
\begin{align}
  s(Q) &= \frac{\text{VWAP}(Q) - P_{\text{mid}}}{P_{\text{mid}}} \times 10{,}000 \quad \text{(buy)} \\
  s(Q) &= \frac{P_{\text{mid}} - \text{VWAP}(Q)}{P_{\text{mid}}} \times 10{,}000 \quad \text{(sell)}
\end{align}
\end{formulabox}

\begin{formulabox}[title={\S9.3 — Maximum Market Order}]
Find the largest quantity that keeps slippage within budget $\beta$:
\begin{equation}
  Q_{\max} = \arg\max_{Q}\; s(Q) \leq \beta
\end{equation}
Uses binary search over the order book levels.
\end{formulabox}


%════════════════════════════════════════════════════════════════
\section{Slippage-at-Risk (SaR)}
%════════════════════════════════════════════════════════════════

\begin{formulabox}[title={\S10.1 — Book Concentration Ratio}]
\begin{equation}
  C = \frac{1}{2}\left(\frac{q_0^{\text{bid}}}{\sum q_i^{\text{bid}}} + \frac{q_0^{\text{ask}}}{\sum q_i^{\text{ask}}}\right)
\end{equation}
$C \to 1$: most liquidity is at level 0 (fragile). $C \to 0$: evenly distributed (robust).
\end{formulabox}

\begin{formulabox}[title={\S10.2 — Concentration Penalty}]
\begin{equation}
  \pi_C = \max\!\left(0,\; 2(C - 0.4)\right) \qquad \text{(additive multiplier)}
\end{equation}
Applied to expected slippage: $\hat{s} = s_{\text{base}} \times (1 + \pi_C)$.
\end{formulabox}

\begin{formulabox}[title={\S10.3 — Slippage-at-Risk}]
\begin{align}
  \text{SaR}_{95\%} &= \hat{s} + 1.645 \cdot \sigma \cdot \sqrt{\frac{T}{24}} \times 10{,}000 \\
  \text{SaR}_{99\%} &= \hat{s} + 2.326 \cdot \sigma \cdot \sqrt{\frac{T}{24}} \times 10{,}000
\end{align}
The VaR-style bounds account for both expected slippage and adverse price movement during execution.
\end{formulabox}

\begin{formulabox}[title={\S10.4 — Book Fragility Signal}]
\begin{equation}
  \text{fragile} = \bigl(C > 0.5\bigr) \;\lor\; \bigl(D_{\text{total}} < 2Q\bigr)
\end{equation}
where $D_{\text{total}} = \sum p_i q_i$ is total depth notional. A fragile book triggers reduced order sizes.
\end{formulabox}


%════════════════════════════════════════════════════════════════
\section{Cascade Detection}
%════════════════════════════════════════════════════════════════

\begin{formulabox}[title={\S11.1 — Spread Ratio}]
\begin{equation}
  R_s = \frac{s_t}{\bar{s}_{\text{rolling}}}
\end{equation}
Current spread divided by rolling mean. $R_s > 3$ signals abnormal widening.
\end{formulabox}

\begin{formulabox}[title={\S11.2 — Volume Ratio}]
\begin{equation}
  R_v = \frac{v_t}{\bar{v}_{\text{rolling}}}
\end{equation}
Current volume rate vs rolling average. $R_v > 5$ with $R_s > 2$ suggests cascade selling/buying.
\end{formulabox}

\begin{formulabox}[title={\S11.3 — Cascade Scoring}]
\begin{equation}
  \text{cascade} = \bigl(R_s > 3\bigr) \;\land\; \bigl(R_v > 5\bigr)
\end{equation}
When detected, execution pauses to avoid adding fuel to the cascade.
\end{formulabox}


%════════════════════════════════════════════════════════════════
\section{Combined Slippage Minimization Objective}
%════════════════════════════════════════════════════════════════

\begin{formulabox}[title={\S12.1 — Global Objective Function}]
The complete execution optimization minimizes:
\begin{equation}
  \boxed{
  \min_{Q_1, \ldots, Q_N} \;\; \sum_{i=1}^{N} \left[
    \underbrace{\eta\sigma\sqrt{\frac{Q_i}{V}}}_{\text{temp.\ impact}}
    + \underbrace{\gamma\frac{Q_i}{V}}_{\text{perm.\ impact}}
    + \underbrace{\lambda\sigma^2\frac{Q_i}{V}\Delta t_i}_{\text{timing risk}}
    + \underbrace{\mathcal{I}_{\text{transient}}(t_i) \cdot \mathbb{1}_{[\text{decay}]}}_{\text{cross-slice}}
  \right]
  }
\end{equation}
subject to:
\begin{align}
  \sum_{i=1}^{N} Q_i &= Q_{\text{total}} & & \text{(full execution)} \\
  Q_i &\leq \rho_{\max} \cdot V_{\text{visible}} & & \text{(participation cap)} \\
  s(Q_i) &\leq \beta_{\max} & & \text{(slippage budget)} \\
  \Delta t_i &\geq \Delta t_{\min} & & \text{(min interval)}
\end{align}
\end{formulabox}

\begin{formulabox}[title={\S12.2 — How Each Module Optimizes This Objective}]
\begin{center}
\begin{tabular}{lll}
\toprule
\textbf{Module} & \textbf{Optimizes} & \textbf{Constraint Enforced} \\
\midrule
Optimal Slicer  & $Q_i$: slippage budget & $s(Q_i) \leq \beta_{\max}$ \\
Flash Limit     & order type selection  & saves spread when $s$ is tight \\
Smart-Market    & $P_{\text{limit}}$: OFI offset & depth-weighted pricing \\
IS Schedule     & $Q_i, \Delta t_i$: Almgren-Chriss & $\min(\text{temp} + \text{timing})$ \\
Transient Model & cross-slice decay     & $\mathcal{I}_{\text{transient}}$ tracking \\
SaR             & risk bounds           & $\text{SaR}_{95\%}$ ceiling \\
Error Handling  & retry/skip/abort      & graceful degradation \\
\bottomrule
\end{tabular}
\end{center}
\end{formulabox}

\end{document}
